% =====================================================================
% 第二讲 (Exposé II) 光滑态射: 一般理论、微分性质
% (Morphismes lisses: généralités, propriétés différentielles)
% A. Grothendieck
% 原书页码: pp. 25-48
% =====================================================================
\chapter{光滑态射：一般理论、微分性质}

对讲 I 的引用以 I 表示. 回忆环都是诺特 (noetherian) 的，预概形都是局部诺特的.

\section{一般理论}

设 $Y$ 是预概形，$t_1, \dots, t_n$ 是未定元，令
\begin{equation}
Y[t_1, \dots, t_n] = Y \otimes_{\mathbf{Z}} \mathbf{Z}[t_1, \dots, t_n]. \tag{1.1}
\label{eq:II.1.1}
\end{equation}
因此 $Y[t_1, \dots, t_n]$ 是 $Y$-概形，在 $Y$ 上仿射 (affine)，由拟凝聚 (quasi-coherent) 代数层 $\mathcal{O}_Y[t_1, \dots, t_n]$ 定义. 该预概形在 $Y$ 上的一个截面的给出，等价于 $\mathcal{O}_Y$ 的 $n$ 个截面的给出 (对应于诸 $t_i$ 在相应同态下的像). 若 $Y'$ 在 $Y$ 之上，则有
\begin{equation}
Y[t_1, \dots, t_n] \times_Y Y' = Y'[t_1, \dots, t_n], \tag{1.2}
\label{eq:II.1.2}
\end{equation}
(这意味着 $Y'$ 到 $Y[t_1, \dots, t_n]$ 的一个 $Y$-态射的给出等价于 $\mathcal{O}_{Y'}$ 的 $n$ 个截面的给出)，另一方面，由 $\mathbf{Z}$ 上多项式环的类似公式，有
\begin{equation}
\bigl(Y[t_1, \dots, t_n]\bigr)[t_{n+1}, \dots, t_m] = Y[t_1, \dots, t_m]. \tag{1.3}
\label{eq:II.1.3}
\end{equation}
公式 \eqref{eq:II.1.2} 蕴含 $Y[t_1, \dots, t_n]$ 随 $Y$ 函子性地变化.

$Y[t_1, \dots, t_n]$ 在 $Y$ 上有限型且平坦 (flat).

\begin{definition}\label{def:II.1.1}
设 $f \colon X \to Y$ 是态射，使 $X$ 成为 $Y$-预概形. 称 $f$ 在 $x \in X$ 处光滑 (smooth)\index[term]{guanghuataishe@光滑态射}\footnote{旧术语：$f$ 在 $x$ 处简单 (simple)，或 $x$ 是 $f$ 的简单点. 这一术语在多种语境 (单代数 (simple algebra)，单群 (simple group)) 中会引起混淆，因而不得不放弃.}，或称 $X$ 在 $Y$ 上于 $x$ 处光滑，若存在整数 $n \geqslant 0$，$x$ 的开邻域 $U$，以及 $U$ 到 $Y[t_1, \dots, t_n]$ 中的平展 $Y$-态射. 称 $f$ (相应地：$X$) 光滑，若它在 $X$ 的所有点处光滑. 环 $A$ 上的代数 $B$ 称为在 $B$ 的素理想 $p$ 处光滑，若 $\mathrm{Spec}(B)$ 在 $\mathrm{Spec}(A)$ 上于 $p$ 点光滑；$B$ 称为在 $A$ 上光滑，若它在 $A$ 上于 $B$ 的每个素理想 $p$ 处光滑. 最后，局部环的局部同态 $A \to B$ 称为光滑的 (或称 $B$ 在 $A$ 上光滑)\footnote{此时最好按 EGA IV 18.6.1 的说法，称 $B$ 在 $A$ 上``本质上光滑'' (essentially smooth)\index[term]{benzhiguanghua@本质光滑}.}，若 $B$ 是 $A$ 上光滑的有限型代数 $B_1$ 的局部化.
\end{definition}

注意 $X$ 在 $Y$ 上的光滑性概念对 $X$ 和 $Y$ 都是局部的；若 $X$ 在 $Y$ 上光滑，则它在 $Y$ 上局部有限型.

% 原书编号笔误: 下面这条命题在原文中与上面的定义同为 ``1.1''
\begingroup
\addtocounter{theorem}{-1}%
\renewcommand{\thetheorem}{\thesection.1}%
\begin{proposition}\label{prop:II.1.1}
$X$ 中 $f$ 光滑的点 $x$ 的集合是开的.
\end{proposition}

\endgroup

由定义即得，平凡.

\begin{corollary}\label{cor:II.1.2}
若 $B$ 在 $A$ 上于 $p$ 处光滑，则它对 $B$ 的每个含于 $p$ 的素理想 $q$ 在 $A$ 上光滑.
\end{corollary}

\ref{prop:II.1.1} 还蕴含 \ref{prop:II.1.1} 中最后两个定义在其共同的存在域内一致.

\begin{proposition}\label{prop:II.1.3}
(i) 平展态射，特别是开浸入、恒等态射，是光滑的. (ii) 光滑态射的基扩张给出光滑态射. (iii) 两个光滑态射的复合是光滑的.
\end{proposition}

\begin{proof}
(i) 由定义即得，平凡，更精确地有：
\end{proof}

\begin{corollary}\label{cor:II.1.4}
平展 = 拟有限 + 光滑.
\end{corollary}

\begin{proof}
(ii) 立即由平展态射 (\ref{prop:II.1.3}) 和投影 $Y[t_1, \dots, t_n] \to Y$ (参见 \eqref{eq:II.1.2}) 的相应事实得出. 对 (iii)，由``平展''(\ref{prop:II.1.3}) 和 $Y[t_1, \dots, t_n] \to Y$ 型投影 (参见 \eqref{eq:II.1.3}) 分别成立，以及 (ii) 中所引的两个事实，形式地得出：设 $Y$ 在 $Z$ 上光滑且 $X$ 在 $Y$ 上光滑，证明 $X$ 在 $Z$ 上光滑；可以假设 $Y$ 在 $Z[t_1, \dots, t_n]$ 上平展且 $X$ 在 $Y[s_1, \dots, s_m]$ 上平展，于是第一个假设蕴含 $Y[s_1, \dots, s_m]$ 在 $Z[t_1, \dots, t_n][s_1, \dots, s_m] = Z[t_1, \dots, s_m]$ 上平展，故 $X$ 在 $Z[t_1, \dots, s_m]$ 上平展，证毕.
\end{proof}

\begin{remark}\label{rem:II.1.5}
定义 \ref{prop:II.1.1} 中的整数 $n$ 是确定好的，因为立即看出它就是 $x$ 的局部环在其纤维 $f^{-1}\bigl(f(x)\bigr)$ 中的维数. 称它为 $X$ 在 $Y$ 上的``相对维数'' (relative dimension)\index[term]{xiangduiweishu@相对维数}. 它对态射的复合具有可加性.
\end{remark}

\section{光滑性的若干判则}

\begin{theorem}\label{thm:II.2.1}
设 $f \colon X \to Y$ 是局部有限型态射，$x \in X$，$y = f(x)$. 为使 $f$ 在 $x$ 处光滑，必要且充分条件是 (a) $f$ 在 $x$ 处平坦，且 (b) $f^{-1}(y)$ 在 $k(y)$ 上于 $x$ 处光滑.
\end{theorem}

\begin{proof}
两个平坦态射的复合平坦，且 $Y[t_1, \dots, t_n] \to Y$ 是平坦态射，故光滑蕴含平坦；结合 \ref{prop:II.1.3} (ii) 即证明必要性. 设 (a) 和 (b) 成立，设 $V$ 是 $y$ 的环为 $A$ 的仿射邻域，$U$ 是 $x$ 在 $V$ 之上的环为 $B$ 的仿射邻域. 取 $U$ 足够小，由 (b) 可以假设存在平展的 $k(y)$-态射
\[
g \colon U|f^{-1}(y) \longrightarrow \mathrm{Spec}\, k[t_1, \dots, t_n] \quad (k = k(y)),
\]
由 $U|f^{-1}(y)$ 的结构层的 $n$ 个截面 $g_i$ 定义. 容易看到可以假设诸 $g_i$ (先验地是 $B \otimes_A k = BS^{-1}$ 的元素，其中 $S = A - \mathfrak{p}$，$\mathfrak{p}$ 是对应于 $y$ 的 $A$ 的素理想) 来自 $U$ 的结构层的截面，因此 $g$ 由一个仍记为 $g$ 的态射
\[
g \colon U \longrightarrow Y[t_1, \dots, t_n]
\]
诱导 (必要时把诸 $g_i$ 同乘以 $k$ 的同一个非零元素). 而由 (a)，$U$ 在 $Y$ 上平坦，$Y[t_1, \dots, t_n]$ 亦然，另一方面 $g$ 在 $y$ 之上的纤维之间诱导平展态射，故由 (\ref{cor:I.5.8})，$g$ 在 $x$ 处平展，证毕.
\end{proof}

\begin{corollary}\label{cor:II.2.2}
设 $S$ 是预概形，$f \colon X \to Y$ 是有限型 $S$-态射，$Y$ 在 $S$ 上有限型且平坦，$x \in X$，$s$ 是 $x$ 在 $S$ 上的投影. 为使 $f$ 在 $x$ 处光滑，必要且充分条件是 $X$ 在 $S$ 上于 $x$ 处平坦 (或：光滑)，且 $f$ 在 $s$ 的纤维上诱导的态射 $f_s \colon X_s \to Y_s$ 在 $x$ 处光滑.
\end{corollary}

\begin{proof}
只需证明充分性，它由判则 \ref{thm:II.2.1} 连同平坦性判则 (\ref{cor:I.5.9}) 得出.
\end{proof}

为了陈述下面的结果，``回忆''一个局部有限型态射 $f \colon X \to Y$ 称为在点 $x \in X$ 处等维 (equidimensional) 的，若 (令 $y = f(x)$) 可以找到 $x$ 的开邻域 $U$，其每个分支都支配 (dominate) $Y$ 的某个分支，使得对每个 $y' \in Y$，$f^{-1}(y') \cap U$ 的不可约分支都有同一个维数，且该维数与 $y'$ 无关. 在这个条件中，只需取 $y'$ 为 $Y$ 的经过 $y$ 的不可约分支的生成点，以及点 $y$. 例如若 $X$ 和 $Y$ 整且 $f$ 支配，则条件意味着经过 $x$ 的 $f^{-1}(y)$ 的分支有``正确''的维数，i.e. 生成纤维的维数 (回忆它们总是 $\geqslant$ 生成纤维的维数). 若 $f$ 在 $x$ 处等维，其纤维在 $x$ 处的维数为 $n$，且 $g \colon U \to Y' = Y[t_1, \dots, t_n]$ 是 $x$ 的邻域 $U$ 的 $Y$-态射，在 $y$ 的纤维上诱导在 $x$ 处拟有限 (或等价地，若 $g$ 在 $x$ 处拟有限) 的态射，则能证明 $U$ 的每个经过 $x$ 的不可约分支都支配 $Y'$ 的某个不可约分支. 此外，由``正规化引理'' (normalization lemma)，这样的 $g$ 总存在 (反过来，若存在 $x$ 的开邻域 $U$ 到形如 $Y' = Y[t_1, \dots, t_n]$ 的 $Y$-概形中的拟有限 $Y$-态射 $g$，使得 $U$ 的每个经过 $x$ 的分支都支配 $Y'$ 的某个分支，则 $f$ 在 $x$ 处等维). 有了这些准备：

\begin{proposition}\label{prop:II.2.3}
设 $f \colon X \to Y$ 是局部有限型态射，$x$ 是 $X$ 的点，$y = f(x)$，假设 $\mathcal{O}_y$ 正规 (normal). 为使 $f$ 在 $x$ 处光滑，必要且充分条件是 $f$ 在 $x$ 处等维，且 $f^{-1}(y)$ 在 $k(y)$ 上于 $x$ 处光滑.
\end{proposition}

\begin{proof}
由定义立即可见光滑态射是等维的\index[term]{dengweitaishe@等维态射} (N.B. 有限型平坦态射未必在 $x$ 处等维，即使其纤维在 $x$ 处不可约). 证明逆命题. 由于 $f^{-1}(y)$ 在 $k(y)$ 上于 $x$ 处光滑，可以假设 (必要时把 $X$ 换成 $x$ 的适当邻域) 存在 $Y$-态射
\[
g \colon X \longrightarrow Y[t_1, \dots, t_n] = Y',
\]
在 $y$ 的纤维上诱导平展态射，从而在 $x$ 处拟有限. 故 $g$ 非分歧，且 ($f$ 在 $x$ 处等维) $X$ 的经过 $x$ 的不可约分支各支配 $Y'$ 的某个分支，更不用说由 $g$ 推出的同态 $\mathcal{O}_{y'} \to \mathcal{O}_x$ (其中 $y' = g(x)$) 单射. 该同态还非分歧，且 $\mathcal{O}_{y'}$ 正规，因为它是环 $\mathcal{O}_y[t_1, \dots, t_n]$ 的局部化，而后者因 $\mathcal{O}_y$ 正规而正规. 故同态 $\mathcal{O}_{y'} \to \mathcal{O}_x$ 平展 (\ref{thm:I.9.5} (ii)).
\end{proof}

\begin{remark}\label{rem:II.2.4}
前面的陈述在把``$\mathcal{O}_y$ 正规''换成更弱的假设``$Y$ 在 $y$ 处几何单支 (geometrically unibranched)''(参见 I 11) 时仍成立 —— 因为 (\ref{thm:I.9.5}) 在该假设下成立. 借此机会同时指出：若整局部环 $A$ 的剩余域代数闭，则解析整 (analytically integral) (i.e. $\widehat{A}$ 整) 蕴含几何单支，反过来在``好环''的任何范畴中都成立，精确地说，在由通常运算封闭、且正规局部环的完备化仍正规的环的范畴中成立 (这一条件，由 Zariski 的``解析正规性定理'' (analytic normality theorem)，在仿射代数及其局部化的范畴中满足)\footnote{参见 EGA IV 7.8.}. 最后在本语境下``回忆''下述 Hironaka 的结果\footnote{参见 EGA IV 5.12.10.}，它有时能保证 $f^{-1}(y)$ 是约化 (reduced) 概形，i.e. 也就是许多代数几何学家滥用为 $f$ 在 $x$ 之上的纤维 (非分歧数) 的东西 (即 $f^{-1}(y)_{\mathrm{r\acute{e}d}}$)：
\end{remark}

\begin{proposition}\label{prop:II.2.5}
设 $f \colon X \to Y$ 是约化预概形的有限型支配态射，$y$ 是 $Y$ 的点，使得 $\mathcal{O}_y$ 正则 (regular). 假设 $f^{-1}(y)$ 的所有分支都重数 (multiplicity) 为 1 (参见下面的定义)，且 $f^{-1}(y)_{\mathrm{r\acute{e}d}}$ 正规. 则 $f^{-1}(y)$ 约化从而正规，$X$ 在 $f^{-1}(y)$ 的所有点处正规，最后 $X$ 在 $f^{-1}(y)$ 的所有点处平坦.
\end{proposition}

称 $f^{-1}(y)$ 的分支 $Z$ 重数为 1，若 $x$ 表示 $Z$ 的生成点，有 (i) $\dim \mathcal{O}_x = \dim \mathcal{O}_y$ (i.e. $Z$ 不是``过剩分支'' (excess component)，即不是``维数太大''的)；(ii) $\mathcal{O}_x$ 的极大理想由 $\mathcal{O}_y$ 的极大理想生成 (后者按 $x$ 的选取先验地生成 $\mathcal{O}_x$ 的一个定义理想).

结合 \ref{prop:II.2.3} 或 \ref{thm:II.2.1}，得到：

\begin{corollary}\label{cor:II.2.6}
设 $f \colon X \to Y$ 是约化预概形的有限型支配态射，$y$ 是 $Y$ 的点，使得 $\mathcal{O}_y$ 正则. 为使 $f$ 在 $X$ 中位于 $y$ 之上的点处光滑，必要且充分条件是 $f^{-1}(y)$ 的分支重数为 1，且 $f^{-1}(y)_{\mathrm{r\acute{e}d}}$ 在 $k(y)$ 上光滑.
\end{corollary}

在过去，当 $Y$ 是离散赋值环 (discrete valuation ring) $A$ 的谱时，人们尤其考虑这种情形，并通常以诸如此类的措辞称之：``若 $X$ 关于所给赋值的约化 (reduction) 是好的''... 此外，那时 $X$ 表示 $\mathbf{P}_K^n$ 的一个闭子概形 ($K$ 是 $A$ 的分式域)，而由于缺乏适当的语言，``定义在 $A$ 上''(而不仅在 $K$ 上) 的对象的更内蕴的作用几乎不显现.

\section{永久性质}

\begin{proposition}\label{prop:II.3.1}
设 $f \colon X \to Y$ 是态射，$x \in X$，$y = f(x)$. 假设 $f$ 在 $x$ 处光滑. 为使 $\mathcal{O}_x$ 约化 (相应地：正则，相应地：正规)，必要且充分条件是 $\mathcal{O}_y$ 约化 (相应地：正则，相应地：正规).
\end{proposition}

\begin{proof}
当 $X$ 形如 $Y[t_1, \dots, t_n]$ 时该陈述确为已知，且对平展态射在 (I，§9) 中已证明；一般情形借助定义 \ref{prop:II.1.1} 立即得出.
\end{proof}

我们不在此详述其他永久性质，它们已经由平坦性本身得出，或由``$X$ 在一个形如 $Y[t_1, \dots, t_n]$ 的 $Y$-预概形之上局部拟有限且平坦''(或者，如我们所说，``$X$ 在 $Y$ 之上 Cohen-Macaulay'') 这一事实得出. 只指出由后一事实得出
\begin{equation}
\dim \mathcal{O}_x = \dim \mathcal{O}_y + n - d, \qquad \operatorname{prof} \mathcal{O}_x = \operatorname{prof} \mathcal{O}_y + n - d, \tag{3.1}
\label{eq:II.3.1}
\end{equation}
其中 $n$ 是 $f$ 的纤维在 $x$ 处的维数，$d$ 是 $k(x)$ 在 $k(y)$ 上的超越次数 (transcendence degree)，于是 (令 coprof = dim $-$ prof)\footnote{对这些公式参见 EGA IV 6.1 和 6.3.}
\begin{equation}
\operatorname{coprof} \mathcal{O}_x = \operatorname{coprof} \mathcal{O}_y. \tag{3.2}
\label{eq:II.3.2}
\end{equation}
例如由此得出 $\mathcal{O}_x$ 是 Cohen-Macaulay 的 (相应地：无嵌入分支) 当且仅当 $\mathcal{O}_y$ 亦然.

\section{光滑态射的微分性质}

为简化起见，我们主要限于一阶微分演算，只对高阶作些快速说明 (那里的结果同样简单).

关于 $Y$-预概形 $X$ 的一次微分 (1-différentielle) 层 $\varOmega_{X/Y}^{1}$ 的定义，参见 (I §1). 假设 $X$ 和 $Y$ 是 $S$-预概形，结构态射 $f \colon X \to Y$ 是 $S$-态射. 则 $f$ 定义一个 (与 $f$ 相容的) 模层 (Module) 同态
\begin{equation}
f^{*} \colon \varOmega_{Y/S}^{1} \longrightarrow \varOmega_{X/S}^{1}, \tag{4.1}
\label{eq:II.4.1}
\end{equation}
换言之，$\varOmega_{X/S}^{1}$ 对 $S$-预概形 $X$ 是反变的. 此外，\eqref{eq:II.4.1} 等价于 $X$ 上的模层同态
\begin{equation}
f^{*}\bigl(\varOmega_{Y/S}^{1}\bigr) \longrightarrow \varOmega_{X/S}^{1}, \tag{4.1bis}
\label{eq:II.4.1bis}
\end{equation}
也姑且记为 $f^{*}$，它嵌入一个模层同态的典范正合序列
\begin{equation}
f^{*}\bigl(\varOmega_{Y/S}^{1}\bigr) \longrightarrow \varOmega_{X/S}^{1} \longrightarrow \varOmega_{X/Y}^{1} \longrightarrow 0. \tag{4.2}
\label{eq:II.4.2}
\end{equation}
所有这些同态都由``局部性质 (归结到仿射情形) 且与算子 $d$ 交换''这一条件定义. \eqref{eq:II.4.2} 的正合性是经典而平凡的，在仿射情形翻译成 (对应于 $A$-代数的同态 $B \to C$) 的正合序列：
\begin{equation}
\varOmega_{B/A}^{1} \otimes_B C \longrightarrow \varOmega_{C/A}^{1} \longrightarrow \varOmega_{C/B}^{1} \longrightarrow 0. \tag{4.2bis}
\label{eq:II.4.2bis}
\end{equation}

\begin{lemma}\label{lem:II.4.1}
设 $f \colon X \to Y$ 是 $S$-预概形的态射. 若 $f$ 非分歧 (相应地：平展)，则 $f^{*}\bigl(\varOmega_{Y/S}^{1}\bigr) \to \varOmega_{X/S}^{1}$ 满射 (相应地：同构). 逆命题在``非分歧''的情形、且 $f$ 假设为局部有限型时成立.
\end{lemma}

\begin{proof}
非分歧的情形由正合序列 \eqref{eq:II.4.2} 和 (\ref{prop:I.3.1}) 得出，但也可以像平展的情形那样直接看. 考虑图
\[
\begin{tikzcd}[row sep=large, column sep=large]
X \arrow[r,"\Delta_{X/Y}"] \arrow[d] & X \times_Y X \arrow[r] & X \times_S X \\
Y \arrow[r,"\Delta_{Y/S}"] & Y \times_S Y
\end{tikzcd}
\]
(原文图式如此)，其中 $X \times_Y X$ 等同于 $Y$ 与 $X \times_S X$ 在 $Y \times_S Y$ 上的纤维积. 由于 $f$ 非分歧，$X \to X \times_Y X$ 是开浸入，故复合浸入 $\Delta_{X/S}$ (由后者与 $X \times_Y X \to X \times_S X$ 复合而得) 的``余法'' (conormal) 层同构于浸入 $X \times_Y X \to X \times_S X$ 的余法层在 $X$ 上的原像. 另一方面，$X \to Y$ 平展从而平坦，故 $X \times_S X \to Y \times_S Y$ 平坦，因此浸入 $X \times_Y X \to X \times_S X$ 的余法层同构于浸入 $Y \to Y \times_S Y$ 的余法层的原像，i.e. $\varOmega_{Y/S}^{1}$ 的原像. 结论由此得出.
\end{proof}

\begin{lemma}\label{lem:II.4.2}
设 $X = Y[t_1, \dots, t_n]$，$Y$ 是 $S$-预概形. 则典范同态序列
\[
0 \longrightarrow f^{*}\bigl(\varOmega_{Y/S}^{1}\bigr) \longrightarrow \varOmega_{X/S}^{1} \longrightarrow \varOmega_{X/Y}^{1} \longrightarrow 0
\]
正合，且 $\varOmega_{X/Y}^{1}$ 以 $d_{X/Y} t_i$ 为基自由.
\end{lemma}

\begin{proof}
(纯仿射的) 验证是直接的. (N.B. 已经知道 \eqref{eq:II.4.2} 的正合性).
\end{proof}

结合这两个陈述和定义 \ref{prop:II.1.1}，得到：

\begin{theorem}\label{thm:II.4.3}
设 $f \colon X \to Y$ 是 $S$-预概形的光滑态射，则：
\begin{enumerate}[(i)]
\item 典范同态序列
\[
0 \longrightarrow f^{*}\bigl(\varOmega_{Y/S}^{1}\bigr) \longrightarrow \varOmega_{X/S}^{1} \longrightarrow \varOmega_{X/Y}^{1} \longrightarrow 0
\]
正合.
\item $\varOmega_{X/Y}^{1}$ 局部自由，其在 $x$ 处的秩 (rank) $n$ 等于 $f$ 在 $x$ 处的相对维数.
\end{enumerate}
\end{theorem}

\begin{corollary}\label{cor:II.4.4}
同态 $f^{*}\bigl(\varOmega_{Y/S}^{1}\bigr) \to \varOmega_{X/S}^{1}$ 单射，其在 $\varOmega_{X/S}^{1}$ 中的像局部是直因子 (direct factor).
\end{corollary}

设 $u \colon F \to G$ 是预概形 $X$ 上的模层同态，称它在 $x \in X$ 处泛单射 (universally injective)，若 $\mathcal{O}_x$-模同态 $F_x \to G_x$ 单射，且与任何 $\mathcal{O}_x$-代数 (或等价地，与任何 $\mathcal{O}_x$-模) 张量后仍单射. 例如，只需存在 $x$ 的开邻域 $U$ 使得 $u$ 诱导 $F|U$ 到 $G|U$ 的某个直因子上的同构；当 $F$ 和 $G$ 在 $x$ 的邻域中自由 (且有限型) 时，这一条件也是必要的，精确地说此时下列条件等价：
\begin{enumerate}[(i)]
\item $u$ 在 $x$ 处单射且 Coker $u$ 在 $x$ 处自由；
\item 存在 $x$ 的开邻域 $U$ 使得 $u$ 诱导 $F|U$ 到 $G|U$ 的某个直因子上的同构；
\item $u$ 在 $x$ 处泛单射；
\item $u$ 在``限制纤维'' $F_x \otimes k(x) \to G_x \otimes k(x)$ 上诱导的同态单射；
\item 转置同态 $\check{G} \to \check{F}$ 在点 $x$ (或等价地，在 $x$ 的邻域) 满射.
\end{enumerate}
(循环证明：(iv)$\Rightarrow$(v) 由 Nakayama 引理得出，另一方面 (v)$\Rightarrow$(i)，因为局部自由商层必为直因子). 几何上，所考虑的情形意味着 $u$ 对应于截面层为 $F$ 的向量丛 (vector bundle) 到由 $G$ 定义的类似向量丛的一个子丛上的同构. 当然，仅凭 $F \to G$ 单射并不够.

\begin{corollary}\label{cor:II.4.5}
设 $f \colon X \to Y$ 是 $S$-预概形的局部有限型态射，$x \in X$，$y = f(x)$，$s$ 是 $x$ 和 $y$ 在 $S$ 上的投影. 假设 $Y$ 在 $y$ 处 $S$-光滑. 下列条件等价：
\begin{enumerate}[(i)]
\item $f$ 在 $x$ 处光滑.
\item $X$ 在 $S$ 上于 $x$ 处光滑，且 $f^{*}\bigl(\varOmega_{Y/S}^{1}\bigr) \to \varOmega_{X/S}^{1}$ 在 $x$ 处泛单射，i.e. 它在 $x$ 处是单射且其余核 $\varOmega_{X/Y}^{1}$ 在 $x$ 处自由.
\end{enumerate}
\end{corollary}

\begin{proof}
必要性由 \ref{prop:II.1.3} (iii) 和 \ref{thm:II.4.3} (i) (ii) 得出，证明充分性. 由于 $dg$ ($g \in \mathcal{O}_x$) 在 $x$ 处生成模 $\varOmega_{X/Y}^{1}$，可以找到 $g_i$ ($1 \leqslant i \leqslant n$) 使得 $dg_i$ 在 $\bigl(\varOmega_{X/Y}^{1}\bigr)_x$ 中的像构成该模的基. 取 $X$ 足够小，可以假设诸 $g_i$ 来自 $\mathcal{O}_X$ 的截面，从而定义 $Y$-态射 $g \colon X \to Y' = Y[t_1, \dots, t_n]$. 利用假设和引理 \ref{lem:II.4.2}，容易看到相应的同态 $g^{*}\bigl(\varOmega_{Y'/S}^{1}\bigr) \to \varOmega_{X/S}^{1}$ 在 $x$ 处双射，这归结为证明：
\end{proof}

\begin{corollary}\label{cor:II.4.6}
设 $f \colon X \to Y$ 是光滑 $S$-预概形的态射. 为使 $f$ 在 $x \in X$ 处平展，必要且充分条件是 $f^{*}\bigl(\varOmega_{Y/S}^{1}\bigr) \to \varOmega_{X/S}^{1}$ 在 $x$ 处是同构.
\end{corollary}

\begin{proof}
由 \ref{lem:II.4.1} 知其必要，且该条件由同一引理蕴含 $f$ 在 $x$ 处非分歧. 由 \ref{cor:II.4.6} 归结到 $S = \mathrm{Spec}(k)$ 的情形. 由于 $Y$ 在 $k$ 上光滑，它正则，从而是正规的，由 (\ref{thm:I.9.5} (ii)) 归结为证明 $\mathcal{O}_y \to \mathcal{O}_x$ 单射，或 $\mathcal{O}_y$ 与 $\mathcal{O}_x$ 维数相同. 而这两个维数分别是 $\varOmega_{Y/k}^{1}$ 和 $\varOmega_{X/k}^{1}$ 在 $y$ 处 (相应地：$x$ 处) 的秩，由假设它们相等.
\end{proof}

\begin{remark}\label{rem:II.4.7}
设 $X$ 和 $Y$ 在 $S$ 上光滑，$f \colon X \to Y$ 的光滑性判则 \ref{cor:II.4.5} (ii) 还可以说成：对每个 $x \in X$，$f$ 在 $x$ 处 (相对基 $S$ 的) 切映射 (tangent map)，i.e. $f^{*}\bigl(\varOmega_{Y/S}^{1}\bigr)$ 与 $\varOmega_{X/S}^{1}$ 在 $x$ 处的限制纤维这两个有限维 $k(x)$-向量空间之间同态的转置，是满射. 这对于使用解析空间 (analytic space) 的人们尤其是一个很熟悉的假设. 他们通常所作的非奇异 (non-singularity) 假设 (它意味着 $X$ 和 $Y$``在 $\mathbf{C}$ 上光滑''，参见 §5) 似乎只是出于代数簇或解析空间的奇异点至今仍使许多几何学家感到恐惧.
\end{remark}

指出 \ref{cor:II.4.6} 的下列特殊情形：

\begin{corollary}\label{cor:II.4.8}
设 $X$ 是 $S$-预概形，$g \colon X \to S[t_1, \dots, t_n]$ 是由 $\mathcal{O}_X$ 的截面 $g_i$ ($1 \leqslant i \leqslant n$) 定义的 $S$-态射，$x$ 是 $X$ 的点，使得 $X$ 在 $S$ 上于 $x$ 处光滑. 为使 $g$ 在 $x$ 处平展，必要且充分条件是 $dg_i$ ($1 \leqslant i \leqslant n$) 在 $x$ 处构成 $\varOmega_{X/S}^{1}$ 的基 (或等价地，它们在 $\varOmega_{X/S}^{1}(x) = \bigl(\varOmega_{X/S}^{1}\bigr)_x \otimes_{\mathcal{O}_x} k(x)$ 中的像构成该 $k(x)$-向量空间的基).
\end{corollary}

\begin{proof}
设 $X$ 是预概形，$Y$ 是由凝聚 (coherent) 理想层 $\mathcal{J}$ 定义的 $X$ 的闭子预概形. 于是 $\mathcal{J}/\mathcal{J}^2$ 可以看作 $Y$ 上的凝聚层 ($Y$ 在 $X$ 中的余法层). 若现在 $X$ 是 $S$-预概形，则有 $Y$ 上的拟凝聚层典范正合序列
\begin{equation}
\mathcal{J}/\mathcal{J}^2 \xrightarrow{d} \varOmega_{X/S}^{1} \otimes_{\mathcal{O}_X} \mathcal{O}_Y \longrightarrow \varOmega_{Y/S}^{1} \longrightarrow 0, \tag{4.3}
\label{eq:II.4.3}
\end{equation}
其右边部分正是 \eqref{eq:II.4.2} (其中 $X$ 与 $Y$ 的角色互换，考虑到 $\varOmega_{Y/X}^{1} = 0$)，而同态 $\mathcal{J}/\mathcal{J}^2 \to \varOmega_{X/S}^{1} \otimes_{\mathcal{O}_X} \mathcal{O}_Y$ 由 (一般非线性) 同态 $g \mapsto dg$ 取商得出. \eqref{eq:II.4.3} 的正合性是经典而平凡的，在仿射情形由下列正合序列解释 (对应于 $A$-代数满同态 $B \to C$，核为 $J$)：
\begin{equation}
J/J^2 \longrightarrow \varOmega_{B/A}^{1} \otimes_B C \longrightarrow \varOmega_{C/A}^{1} \longrightarrow 0 \quad (C = B/J). \tag{4.3bis}
\label{eq:II.4.3bis}
\end{equation}
(这个正合序列在 (\ref{prop:II.4.9}) 的证明中其实已被隐含使用!)
\end{proof}

\begin{proposition}\label{prop:II.4.9}
设 $X$ 是 $S$-预概形，$Y$ 是由 $X$ 上的凝聚理想层 $\mathcal{J}$ 定义的 $X$ 的闭子预概形，$x$ 是 $X$ 的点，$g_i$ ($1 \leqslant i \leqslant n$) 是 $\mathcal{O}_X$ 的截面，定义 $S$-态射
\[
g \colon X \longrightarrow S[t_1, \dots, t_n] = X',
\]
最后 $p$ 是整数，$0 \leqslant p \leqslant n$. 假设 $X$ 在 $S$ 上于 $x$ 处光滑. 下列条件等价：
\begin{enumerate}[(i)]
\item 存在 $x$ 在 $X$ 中的开邻域 $X_1$，使得 $g|X_1$ 平展，且 $Y_1 = Y \cap X_1$ ($Y$ 在 $X_1$ 上的迹) 是 $X' = S[t_1, \dots, t_n]$ 的闭子预概形 $Y' = S[t_{p+1}, \dots, t_n]$ 的原像 (i.e. $g_i$ ($1 \leqslant i \leqslant p$) 生成 $\mathcal{J}|X_1$)；
\item $Y$ 在 $S$ 上于 $x$ 处光滑，$g_i$ ($1 \leqslant i \leqslant p$) 定义 $\mathcal{J}_x$ 的元素，$dg_i(x)$ ($1 \leqslant i \leqslant n$) 构成 $k(x)$ 上的基 $\varOmega_{X/S}^{1}(x)$，$dg'_i(x)$ ($p+1 \leqslant i \leqslant n$) 构成 $k(x)$ 上的基 $\varOmega_{Y/S}^{1}(x)$ (其中 $g'_i$ 表示 $g_i$ 到 $Y$ 的限制；微分是关于 $S$ 取的)；
\item $g_i$ ($1 \leqslant i \leqslant p$) 定义 $\mathcal{J}_x$ 的生成元组，且 $dg_i(x)$ ($1 \leqslant i \leqslant n$) 构成 $k(x)$ 上的基 $\varOmega_{X/S}^{1}(x)$；
\item $Y$ 在 $S$ 上于 $x$ 处光滑，$g_i$ ($1 \leqslant i \leqslant p$) 构成 $\mathcal{J}_x$ 的极小生成元组，$dg'_i(x)$ ($p+1 \leqslant i \leqslant n$) 构成 $k(x)$ 上的基 $\varOmega_{Y/S}^{1}(x)$.
\end{enumerate}
此外，在这些条件下，$\mathcal{J}/\mathcal{J}^2$ 是 $Y$ 上在 $x$ 处自由的模层，以诸 $g_i$ ($1 \leqslant i \leqslant p$) 的类为 $x$ 处的基，且典范同态 $\mathcal{J}/\mathcal{J}^2 \to \varOmega_{X/S}^{1} \otimes \mathcal{O}_Y$ 在 $x$ 处泛单射.
\end{proposition}

\begin{remarkstar}
这意味着 $p$ 由其他条件确定，或作为 $\mathcal{J}/\mathcal{J}^2$ 在 $Y$ 上于 $x$ 处的自由模层的秩，或作为 $\mathcal{J}_x$ 在 $X$ 上的生成元的最小个数，或最后由``$Y$ 相对 $S$ 在 $x$ 处的相对维数为 $n - p$''这一事实确定.
\end{remarkstar}

\paragraph*{证明.} 先假设 (i) 成立. 则由 (\ref{prop:II.4.9} (iii))，$Y_1$ 在 $Y'$ 上平展，故按定义它在 $S$ 上于 $x$ 处光滑 (相对维数 $n - p$)，$Y$ 亦然. 于是由 (\cref{cor:II.4.8}) 得出 $dg_i$ ($1 \leqslant i \leqslant n$) 构成 $\varOmega_{X/S}^{1}$ 在 $x$ 处的基，$dg'_i$ ($p+1 \leqslant i \leqslant n$) 构成 $\varOmega_{Y/S}^{1}$ 在 $x$ 处的基，由此由正合序列 \eqref{eq:II.4.3} 得出 $g_i$ ($1 \leqslant i \leqslant p$) 在 $\mathcal{J}/\mathcal{J}^2$ (看作 $Y$ 上的模层) 中于 $x$ 处线性无关；由于 $g_i$ ($1 \leqslant i \leqslant p$) 生成 $\mathcal{J}_x$，故 $g_i$ 模 $\mathcal{J}_x^2$ 构成 $\mathcal{J}/\mathcal{J}^2$ 在 $x$ 处的基. 这既蕴含 $g_i$ ($1 \leqslant i \leqslant p$) 构成 $\mathcal{J}_x$ 的极小生成元组，又蕴含 \eqref{eq:II.4.3} 的同态 $\mathcal{J}/\mathcal{J}^2 \to \varOmega_{X/S}^{1} \otimes \mathcal{O}_Y$ 在 $x$ 处泛单射 (因为它把一个 $x$ 处自由模层的基映到 $x$ 处自由模层的基的一部分 —— N.B. 这是 $Y$-模层). 这证明 (i) 蕴含 (ii)，(iii)，(iv) 以及命题 \ref{prop:II.4.9} 的最后断言.

(iii) 由推论 \ref{cor:II.4.8} 蕴含 (i).

(ii) 蕴含 (i). 事实上，(ii) 中的第一个假设意味着 (把 $X$ 换成 $x$ 在 $X$ 中的开邻域) $g$ 诱导态射 $h \colon Y \to Y'$. 由 \ref{cor:II.4.8}，(ii) 的其他两个假设意味着 $g$ 在 $x$ 处平展，$h$ 在 $x$ 处平展. 设 $Y''$ 是 $Y'$ 在 $g$ 下的原像. 于是 $Y$ 是 $Y''$ 的闭子预概形，而由 (\ref{prop:II.4.9} (iii))，$Y''$ 在 $Y'$ 上于 $x$ 处平展 (因为 $g$ 在 $x$ 处平展). 故浸入态射 $Y \to Y''$ 本身平展 (\ref{cor:II.4.8})，从而是开浸入 (\ref{cor:I.5.8} 或 \ref{cor:I.5.2})，于是再把 $X$ 换成 $x$ 的适当开邻域 $X_1$，得到 (i).

以上建立了条件 (i)，(ii)，(iii) 的等价性，以及它们蕴含 (iv)；剩下证明 (iv)$\Rightarrow$(ii)，一旦知道``$Y$ 在 $S$ 上于 $x$ 处光滑蕴含 $\mathcal{J}/\mathcal{J}^2$ 在 $Y$ 上于 $x$ 处自由，且同态 $\mathcal{J}/\mathcal{J}^2 \to \varOmega_{X/S}^{1} \otimes \mathcal{O}_Y$ 在 $x$ 处泛单射''(考虑到 $\varOmega_{X/S}^{1}$ 在 $X$ 上于 $x$ 处自由) 就立即得出. 最后这一点包含在

\begin{theorem}\label{thm:II.4.10}
设 $X$ 是光滑 $S$-预概形，$Y$ 是由 $X$ 上的凝聚理想层 $\mathcal{J}$ 定义的 $X$ 的闭子预概形，$x$ 是 $X$ 的点. 下列条件等价：
\begin{enumerate}[(i)]
\item $Y$ 在 $S$ 上于 $x$ 处光滑.
\item 存在 $x$ 在 $X$ 中的开邻域 $X_1$ 和平展 $S$-态射
\[
g \colon X_1 \longrightarrow X' = S[t_1, \dots, t_n],
\]
使得 $Y_1 = Y \cap X_1$ ($Y$ 在 $X_1$ 上的迹) 是 $X_1$ 中由 $g$ 从 $X' = S[t_1, \dots, t_n]$ 的闭子预概形 $Y' = S[t_{p+1}, \dots, t_n]$ 取原像得到的子预概形，对某个适当的 $p$.
\item 存在 $\mathcal{J}_x$ 的生成元 $g_i$ ($1 \leqslant i \leqslant p$)，使得 $dg_i$ 构成 $\varOmega_{X/S}^{1}$ 在 $x$ 处的基的一部分 (或等价地，使 $dg_i(x)$ 在 $\varOmega_{X/S}^{1}$ 中在 $k(x)$ 上线性无关).
\item 层 $\mathcal{J}/\mathcal{J}^2$ 在 $Y$ 上于 $x$ 处自由，且典范同态
\[
d \colon \mathcal{J}/\mathcal{J}^2 \longrightarrow \varOmega_{X/S}^{1} \otimes \mathcal{O}_Y
\]
在 $x$ 处泛单射；或：典范同态序列
\[
0 \longrightarrow \mathcal{J}/\mathcal{J}^2 \longrightarrow \varOmega_{X/S}^{1} \otimes \mathcal{O}_Y \longrightarrow \varOmega_{Y/S}^{1} \longrightarrow 0
\]
在 $x$ 处正合，且 $\varOmega_{Y/S}^{1}$ 在 $x$ 处局部自由.
\end{enumerate}
\end{theorem}

\paragraph*{证明.} 由前面的讨论已知 (ii) 蕴含 (i)，(iii)，(iv). 证明 (i)$\Rightarrow$(ii) (这将同时完成 \ref{prop:II.4.9} 的证明). 由定理 \ref{thm:II.4.10} (ii)，正合序列 \eqref{eq:II.4.3} 中最后两项是 $Y$ 上的自由模层. 于是，由于 $dg$ ($g \in \mathcal{O}_X$) 在 $\varOmega_{X/S}^{1} \otimes_{\mathcal{O}_X} \mathcal{O}_Y$ 中的像在 $x$ 处生成该模，故它们在 $\varOmega_{Y/S}^{1}$ 中的像在 $x$ 处生成后者，可以找到 $g_i$ ($p+1 \leqslant i \leqslant n$) $\in \mathcal{O}_X$ 使得 $dg'_i$ 构成 $\varOmega_{Y/S}^{1}$ 的基，然后 (由 \eqref{eq:II.4.3} 的正合性) 用形如 $dg_i$ ($1 \leqslant i \leqslant p$，诸 $g_i \in \mathcal{J}_x$) 的元素把 $dg_i$ ($p+1 \leqslant i \leqslant n$) 的组补成中项的基. 诸 $g_i$ 来自 $X$ 中 $x$ 的邻域上的 $\mathcal{O}_X$ 的截面，可以假设该邻域等于 $X$. 于是处于 \ref{cor:II.4.8} (ii) 的条件下，而我们已经建立了这蕴含条件 \ref{cor:II.4.8} (i)，由此得 \ref{thm:II.4.10} (ii).

\ref{thm:II.4.10} 中 (iii)$\Rightarrow$(ii) 由 \ref{cor:II.4.8} 中 (iii)$\Rightarrow$(i) 立即得出. 故 (i)，(ii)，(iii) 等价，且蕴含 (iv). 最后，(iv) 蕴含 (iii) 是平凡的，考虑到模 $\mathcal{J}_x^2$ 构成 $\mathcal{J}_x$ 的基的 $g_i \in \mathcal{J}_x$ 生成 $\mathcal{J}_x$ (Nakayama 引理).

此外，前面的证明还显示：

\begin{corollary}\label{cor:II.4.11}
设 $X$ 是 $S$-预概形，$Y$ 是由 $X$ 上的凝聚理想层 $\mathcal{J}$ 定义的闭子预概形，$x$ 是 $Y$ 的点. 假设 $X$ 和 $Y$ 在 $S$ 上于 $x$ 处光滑. 设 $g_i$ 是 $\mathcal{J}$ 的截面 ($1 \leqslant i \leqslant p$). 下列条件等价：
\begin{enumerate}[(i)]
\item $g_i$ 生成 $\mathcal{J}_x$，且 $dg_i(x)$ 在 $\varOmega_{X/S}^{1}(x)$ 中在 $k(x)$ 上线性无关；
\item $g_i$ 模 $\mathcal{J}^2$ 构成 $\mathcal{J}/\mathcal{J}^2$ 在 $x$ 处的基；
\item $g_i$ 构成 $\mathcal{J}_x$ 的极小生成元组；
\item 可以找到 $X$ 的邻域 $X_1$ 上 $\mathcal{O}_X$ 的另外一些截面 $g_i$ ($p+1 \leqslant i \leqslant n$)，与前面的截面一起定义平展态射 $X_1 \to X' = S[t_1, \dots, t_n]$，使得 $Y_1 = Y \cap X_1$ 是 $g$ 从 $X' = S[t_1, \dots, t_n]$ 的闭子预概形 $Y' = S[t_{p+1}, \dots, t_n]$ 取原像得到的.
\end{enumerate}
\end{corollary}

特别地：

\begin{corollary}\label{cor:II.4.12}
设 $X$ 是 $S$-预概形，$F$ 是 $\mathcal{O}_X$ 的截面，$Y$ 是 $F$ 的零点子概形 (由凝聚理想 $F \cdot \mathcal{O}_X$ 定义)，$x$ 是 $Y$ 的点. 假设 $X$ 在 $S$ 上于 $x$ 处简单. 为使 $Y$ 在 $S$ 上于 $x$ 处光滑，必要且充分条件是 $F$ 在 $x$ 的邻域内为零，或 $dF(x) \neq 0$ (其中 $dF(x)$ 表示 $dF$ 在 $k(x)$ 上的向量空间 $\varOmega_{X/S}^{1}(x)$ 中的像).
\end{corollary}

\begin{proof}
由 \ref{thm:II.4.10} 判则 (iii) 知其充分. 它也是必要的，因为 $\mathcal{J}$ 由一个元素生成，首先必须 $\mathcal{J}/\mathcal{J}^2$ 在点 $x$ 处自由且秩 $\leqslant 1$. 若此秩为 0，i.e. $\mathcal{J}/\mathcal{J}^2 = 0$ 在 $x$ 处，则由 Nakayama 引理推出 $\mathcal{J} = 0$ 在 $x$ 处，i.e. $F$ 在 $x$ 的邻域内为零. 若此秩为 1，则 $F$ 构成 $\mathcal{J}$ 在 $x$ 处的极小生成元组，由 (\cref{cor:II.4.11}) 中 (i) 与 (iii) 的等价性即得结论.
\end{proof}

\begin{corollary}\label{cor:II.4.13}
设 $Y$ 是局部有限型 $S$-预概形，$S'$ 是平坦 $S$-预概形，$Y' = Y \times_S S'$，$x'$ 是 $Y'$ 的点，$x$ 是它在 $Y$ 中的典范像. 为使 $Y$ 在 $S$ 上于 $x$ 处光滑，必要且充分条件是 $Y'$ 在 $S'$ 上于 $x'$ 处光滑. 特别地，若 $S' \to S$ 平坦且满射，则 $Y$ 在 $S$ 上光滑当且仅当 $Y'$ 在 $S'$ 上光滑.
\end{corollary}

\begin{proof}
只需证明充分性 (必要性已在 \ref{prop:II.1.3} (ii) 中看到). 可以假设 (把 $Y$ 换成 $x$ 的适当邻域，$Y'$ 换成它的原像) $Y$ 是仿射 $S$ 上的仿射有限型概形，故 $Y$ 同构于 $S[t_1, \dots, t_n]$ 的闭子概形. 于是 $Y'$ 等同于 $X' = X \times_S S'$ 的闭子预概形. 由于 $X$ 在 $S$ 上光滑，$X'$ 在 $S'$ 上光滑，可以应用光滑性判则 \ref{thm:II.4.10}. 这里判则 (iv) 立即给出结果.
\end{proof}

\begin{remark}\label{rem:II.4.14}
定理 \ref{thm:II.4.10} 的判则 (iii) 值得称为 \emph{Jacobi 光滑性判则} (Jacobian criterion of smoothness)\index[term]{Jacobi guanghuaxingpanze@Jacobi 光滑性判则}. 它使得理论上能够识别一个给定的 $S$-预概形 $Y$ 是否在 $S$ 上于 $Y$ 的点 $x$ 处光滑，因为 $Y$ 的邻域总同构于光滑 $S$-预概形 $X$ 的子概形，例如 $X = S[t_1, \dots, t_n]$. 事实上，通常正是在 $X = S[t_1, \dots, t_n]$，$S = \mathrm{Spec}(A)$ 时陈述 Jacobi 判则 (当然，在 Zariski 所考虑的经典情形，$A$ 是域). 把与 $A[t_1, \dots, t_n]$ 的理想 $J$ 及包含它的素理想有关的陈述留给读者. 注意到目前看来 (尤其自从 Nagata 用非微分方法推广了 Zariski 的``代数概形的正则点集是开的''这一定理以来)，Jacobi 判则几乎只有以我们这里所给出的形式才有意义 (i.e. 只使用相对微分，而不使用绝对微分，i.e. 相对绝对常数环 $\mathbf{Z}$ 的微分). 和常常见到的情况一样，在这里考虑微分比考虑导子更方便. 最后注意，若 $Y$ 在 $S$ 上于 $x$ 处光滑，相对维数为 $n$，则存在 $Y$ 中 $x$ 的开邻域同构于 $X = S[t_1, \dots, t_n]$ 的子概形，其中 $n = m + 1$，这由定义和 (\ref{thm:II.4.10}) 得出.
\end{remark}

设 $A$ 是诺特环，$x_i$ ($1 \leqslant i \leqslant n$) 是 $A$ 的元素，$J$ 是由诸 $x_i$ 生成的理想. 称诸 $x_i$ 构成 $J$ 的正则生成元组 (regular system of generators)，若典范满同态
\[
(A/J)[t_1, \dots, t_n] \longrightarrow \mathrm{gr}^{J}(A)
\]
(由诸 $x_i$ 定义，其中右端表示与 $A$ 按 $J$ 的幂滤链所得的相伴分次环) 是同构. 这一条件等价于下列两个条件：
\begin{enumerate}[(i)]
\item 典范满同态
\[
S_{A/J}(J/J^2) \longrightarrow \mathrm{gr}^{J}(A)
\]
(其中左端表示 $A/J$-模 $J/J^2$ 的对称代数) 是同构，且
\item $J/J^2$ 自由，以诸 $x_i$ 模 $J^2$ 的类为基.
\end{enumerate}
以这种形式看出，若 $J \neq A$，则诸 $x_i$ 构成 $J$ 的极小生成元组，且 $J$ 的任何其他极小生成元组都是正则生成元组 (N.B. ``极小'' 取严格意义：元素个数最少，它仅当 $A$ 局部时才等价于``按包含关系极小''的意义)；另一方面，若 $J = A$，则 $J$ 的任何生成元组都是正则的.

理想的正则生成元组的正则性对任意乘法稳定集 (multiplicatively closed set) 的局部化稳定，另一方面立即看出，为使 $(x_i)$ 是 $J$ 的极小生成元组，只需对包含 $J$ 的每个极大理想 $m$，$x_i$ 在 $A_{\mathfrak{m}}$ 中定义 $JA_{\mathfrak{m}}$ 的正则生成元组. 因此归结到 $A$ 是以 $m$ 为极大理想的局部环、诸 $x_i$ 属于 $m$ 的情形. 此时诸 $x_i$ 构成 $J$ 的正则生成元组，当且仅当它们构成 Serre 意义下的 $A$-序列 (A-sequence)\footnote{现在我们宁愿说 ``$A$-正则序列'' (A-regular sequence)，参见 EGA 0$_{\mathrm{IV}}$ 15.1.7 和 15.1.11.}，i.e. 若对每个满足 $1 \leqslant i \leqslant n$ 的 $i$，$x_i$ 在 $A/(x_1, \dots, x_{i-1})A$ 中非零因子 (non-diviseur de 0).

最后，当 $A$ 是环 $B$ 上的代数，且 $A/J$ 作为 $B$-代数同构于 $B$ (从而 $J$ 是 $B$-代数同态 $A \to B$ 的核) 时，诸 $x_i$ 构成 $J$ 的正则生成元组，当且仅当典范同态
\[
B[\![t_1, \dots, t_n]\!] \longrightarrow \widehat{A}
\]
(由诸 $x_i$ 定义，其中右端表示 $A$ 按 $J$ 的幂所定义的拓扑的分离完备化 $\varprojlim A/J^{n+1}$) 是同构 (它总是满射).

所有这些事实都是众所周知的，大概几乎完全见于 Serre 的交换代数课程 (由 Gabriel 整理) 中 (那里，在 $A$ 是局部环的情形，还可以找到 $A$-序列的其他刻画).

设 $J$ 是诺特环 $A$ 中的理想. 称 $J$ 为\emph{正则理想} (regular ideal)，若对 $A$ 的每个素理想 $\mathfrak{p}$，$JA_{\mathfrak{p}}$ 有正则生成元组. 显然只需对 $\mathfrak{p} \supset J$ 验证，且可以只限于 $p$ 极大. 更一般地，设 $\mathcal{J}$ 是局部诺特预概形 $X$ 上的理想层，称 $\mathcal{J}$ 为正则理想层，若对每个 $x \in X$，$\mathcal{J}_x$ 是 $\mathcal{O}_x$ 的有正则生成元组的理想. 这等价于下列两个条件的合取：
\begin{enumerate}[a)]
\item 典范满同态
\[
S_{\mathcal{O}_X/\mathcal{J}}(\mathcal{J}/\mathcal{J}^2) \longrightarrow \mathrm{gr}^{\mathcal{J}}(\mathcal{O}_X)
\]
是同构，且
\item $\mathcal{O}_X/\mathcal{J}$-模层 $\mathcal{J}/\mathcal{J}^2$ 局部自由.
\end{enumerate}
此时还说由 $\mathcal{J}$ 定义的 $X$ 的子预概形 $Y$ (即 $\mathcal{O}_Y$ 以 0 延拓后同构于 $\mathcal{O}_X/\mathcal{J}$ 者) 在 $X$ 中正则浸入 (regularly immersed)，并同样 (以显然的方式) 定义正则浸入态射 (morphisme d'immersion régulière) (相应地：在点 $x$ 处正则) 的概念，即把 $Y$ (相应地：$x$ 的适当邻域) 等同于 $X$ 的某个开集中的正则浸入闭子预概形的浸入态射. (不应当说：正则子预概形，因为那将意味着 $Y$ 的局部环正则). 最后，$\mathcal{J}$ 的截面 $x_i$ 称为正则生成元组，若对每个 $x \in X$，$\mathcal{O}_x$ 中相应元素构成 $\mathcal{J}_x$ 的正则生成元组 (该术语与为环的理想生成元引入的术语相容). 这也意味着由诸 $x_i$ 定义的典范满同态
\[
\mathcal{O}_Y[t_1, \dots, t_n] \longrightarrow \mathrm{gr}^{\mathcal{J}}(\mathcal{O}_X)
\]
是同构. 若事先知道理想层 $\mathcal{J}$ 是正则的\index[term]{zhengzelixiang@正则理想}，这也简单意味着在 $Y$ 的每个点 $x$ 处，诸 $x_i$ 定义 $\mathcal{J}/\mathcal{J}^2$ 在 $\mathcal{O}_{Y,x}$ 上的基. (N.B. 若 $Y$ 为空，该条件为空.) 于是，为使 $\mathcal{J}$ 有正则生成元组\index[term]{zhengzeshengchengyuanxi@正则生成元系}，必要且充分条件是 $\mathcal{J}$ 正则，且 $\mathcal{O}_Y$-模层 $\mathcal{J}/\mathcal{J}^2$ 整体自由 (而不只是局部自由)，i.e. 典范同态 $S_{\mathcal{O}_Y}(\mathcal{J}/\mathcal{J}^2) \to \mathrm{gr}^{\mathcal{J}}(\mathcal{O}_X)$ 满射，且 $\mathcal{O}_Y$-模层 $\mathcal{J}/\mathcal{J}^2$ 整体自由.

增广环 (augmented ring) 称为正则的\index[term]{zhengzezengguanghuan@正则增广环}，若增广理想 (idéal d'augmentation) 正则. 于是，若 $A$ 是局部环，看作增广到其剩余域 $k$，则 $A$ 是正则局部环当且仅当它是\emph{正则增广环}.

(老实说，似乎本不必先从环开始做这套初步推论 (sorites)，而最好立即从层开始. 若想要诺特情形的某种东西，这里所采纳的定义 —— 先验地比 Serre 的 $A$-序列定义更宽松 —— 似乎对微分演算的需要更可取. 当然，要做得好的话，还应当在非诺特框架下至少发展一部分光滑态射理论\footnote{如前言所说，这现在已实现，参见 EGA IV 17，18.}，大概从 Jacobi 判则出发，以便尽可能得到光滑态射和\emph{平展态射} (i.e. 光滑且拟有限) 的一切本质形式性质；只有逆命题才用到诺特假设.)

在这些冗长的术语准备之后，一个小定理：

\begin{theorem}\label{thm:II.4.15}
设 $X$ 是局部有限型 $S$-预概形，$Y$ 是由 $X$ 上的凝聚理想层 $\mathcal{J}$ 定义的 $X$ 的闭子预概形，$x$ 是 $X$ 的点. 现在假设 $Y$ 在 $S$ 上于 $x$ 处光滑 (对 $X$ 不作假设). 则下列条件等价：
\begin{enumerate}[(i)]
\item $X$ 在 $S$ 上于 $x$ 处光滑；
\item 浸入 $i \colon Y \to X$ 在 $x$ 处正则，i.e. $\mathcal{J}_x$ 是 $\mathcal{O}_x$ 的正则理想.
\end{enumerate}
\end{theorem}

\begin{corollary}\label{cor:II.4.16}
假设 $Y$ 在 $S$ 上光滑. 为使 $X$ 在 $Y$ 的邻域中 (i.e. 在 $Y$ 的点处) 在 $S$ 上光滑，必要且充分条件是 $Y$ 正则嵌入 $X$，i.e. 浸入 $i \colon Y \to X$ 正则.
\end{corollary}

\paragraph*{证明.} (i) 蕴含 (ii). 应用 \ref{thm:II.4.10} 判则 (ii)，由于 $g \colon X_1 \to X$ 平坦，为证明 $X'$ 的子预概形 $Y'$ 在 $g$ 下的原像正则嵌入，归结为证明 $Y' = S[t_{p+1}, \dots, t_n]$ 正则嵌入 $S[t_1, \dots, t_n]$，这是平凡的 ($t_i$ ($1 \leqslant i \leqslant p$) 构成定义 $X'$ 中 $Y'$ 的理想的正则生成元组).

(ii) 蕴含 (i). 设 $g_i$ ($1 \leqslant i \leqslant p$) 是 $\mathcal{J}_x$ 的正则生成元组，$g_i$ ($p+1 \leqslant i \leqslant n$) 是 $\mathcal{O}_{X,x}$ 的元素，使得它们在 $\mathcal{O}_{Y,x}$ 中的像 $g'_i$ 定义平展态射
\[
Y_1 \longrightarrow Y' = S[t_{p+1}, \dots, t_n]
\]
(把 $Y$ 的邻域 $Y_1$ 映到 $Y'$). 诸 $g_i$ ($1 \leqslant i \leqslant n$) 来自 $x$ 的邻域 $X_1$ 上 $\mathcal{O}_X$ 的 (同名) 截面，且可以假设 $X_1 = X$，$Y_1 = Y$. 由此得到态射
\[
g \colon X \longrightarrow X' = S[t_1, \dots, t_n],
\]
一切归结为证明该态射在 $x$ 处平展. 取 $X_1$ 足够小，可以假设 $g_i$ ($1 \leqslant i \leqslant p$) 在整个 $X$ 上构成 $\mathcal{J}$ 的正则生成元组. 特别地，它们生成 $\mathcal{J}$，故 $X$ 的子预概形 $Y$ 等同于 $g$ 从 $X'$ 的子预概形 $Y'$ 取原像得到的子预概形. 设 $x' = g(x)$，则 $X' \to X$ 在 $x'$ 处的纤维等同于 $Y \to Y'$ 在 $x$ 处的纤维，故在 $k(x')$ 上平展，于是 $g$ 在 $x$ 处非分歧，剩下证明 $g$ 在 $x$ 处平坦. 而按 $\mathcal{J}'_x$ 的幂滤链所得的 $\mathcal{O}_{X',x'}$ 的相伴分次环在各次数都是 $\mathcal{O}_{Y',x'}$ 上的自由环，另一方面按 $\mathcal{J}_x = \mathcal{J}'_x \mathcal{O}_{X,x}$ 的幂滤链所得的 $\mathcal{O}_{X,x}$ 的相伴分次环 (在典范同态下) 同构于前者的 $\mathcal{O}_{Y,x}$ 张量 (因为两个环都是常数环 $\mathcal{O}_{Y',x'}$ (相应地：$\mathcal{O}_{Y,x}$) 上 $n - p$ 个未定元的多项式环)，最后在 $\mathcal{O}_{X',x'}/\mathcal{J}'_{x'} = \mathcal{O}_{Y',x'}$ 上 $\mathcal{O}_{X,x}/\mathcal{J}_x = \mathcal{O}_{Y,x}$ 平坦.

由平坦性的一般判则 (适用于诺特局部环的局部同态 $A' \to A$，$A'$ 带有理想 $J' \neq A'$，使得相伴分次环在 $A'/J'$ 上各维数自由) 推出 $X$ 在 $X'$ 上于 $x$ 处平坦，证毕.

\begin{corollary}\label{cor:II.4.17}
设 $X$ 是在 $Y$ 上局部有限型的预概形，$i$ 是 $X$ 在 $Y$ 上的截面，$y$ 是 $Y$ 的点，$x = i(y)$，$\mathcal{J}$ 是由子预概形 $i(Y)$ 定义的 $X$ 上的理想层 (为简化陈述我们假设它闭，当 $X$ 是概形时此条件成立).

下列条件等价：
\begin{enumerate}[(i)]
\item $X$ 在 $Y$ 上于 $x$ 处光滑；
\item $i$ 在 $y$ 处是正则浸入；
\item $\mathcal{O}_x$ 按 $\mathcal{J}_x$ 的幂所定义的拓扑的完备化 $\mathcal{O}_y$-代数是形式幂级数代数 $\mathcal{O}_y[\![t_1, \dots, t_n]\!]$；
\item[(iii bis)] 存在 $y$ 的开邻域 $U$，使得 $\mathcal{O}_Y$ 上的代数层 $\varprojlim i^{*}\bigl(\mathcal{O}_X/\mathcal{J}^{n+1}\bigr)$ 同构于 $U$ 之上的形如 $\mathcal{O}_Y[\![t_1, \dots, t_n]\!]$ 的层；
\item 存在 $y$ 的开邻域 $U$，$x$ 的开邻域 $V$，以及 $Y$-态射 $g \colon V \to U[t_1, \dots, t_n]$，使得 $g$ 平展，$i$ 诱导 $V$ 在 $U$ 上的截面，经 $g$ 变为 $U[t_1, \dots, t_n]$ 在 $U$ 上的零截面.
\end{enumerate}
\end{corollary}

\begin{proof}
(i) 与 (ii) 的等价性是定理 \ref{thm:II.4.15} 在 $Y = S$ 时的特殊情形；(ii) 与 (iii) 的等价性 (以及实质上 (ii) 与 (iii bis) 的等价性) 已在``回忆''中指出；至于 (i) 与 (iv) 的等价性，容易由定理 \ref{thm:II.4.15} (其条件 (i) 与 (ii) 的等价性) 推出.
\end{proof}

\begin{corollary}\label{cor:II.4.18}
设 $X$ 是在 $S$ 上光滑的预概形. 则对角态射 (diagonal morphism)
\[
\Delta_{X/S} \colon X \longrightarrow X \times_S X
\]
是正则浸入，或如人们所说的，$X$ 在 $S$ 上``微分光滑''\index[term]{weifenguanghua@微分光滑} (differentially smooth).
\end{corollary}

\begin{proof}
事实上，这是推论 \ref{cor:II.4.16} 的特殊情形，因为 $X$ 和 $X \times_S X$ 都在 $S$ 上光滑.
\end{proof}

% 原书编号笔误: 下面这条在原文中与上面的推论同为 ``\ref{cor:II.4.18}''
\begingroup
\addtocounter{theorem}{-1}%
\renewcommand{\thetheorem}{\thesection.18}%
\begin{remark}\label{rem:II.4.18bis}
回忆 (I 1)：若 $X$ 是 $S$ 上的预概形，引入 $X$ 上的拟凝聚代数层 $\mathcal{P}_{X/S}^{n} = \mathcal{O}_{X \times_S X}/\mathcal{I}_X^{n+1}$ (其中 $\mathcal{I}_X$ 表示在 $X \times_S X$ 中定义对角线的理想层)，借助第一投影 $\mathrm{pr}_1 \colon X \times_S X \to X$ 把它看作 $\mathcal{O}_X$-代数层. 诸 $\mathcal{P}_{X/S}^{n}$ 构成 $X$ 上的代数层投射系 (projective system)，其投射极限记为 $\hat{\mathcal{P}}_{X/S}^{\infty}$，它就是 $X \times_S X$ 沿对角线的形式完备化 (formal completion) 的结构层 (现在假设 $X$ 在 $S$ 上局部有限型，故诸 $\mathcal{P}_{X/S}^{n}$ 凝聚). 说 $X$ 在 $S$ 上\emph{微分光滑} (i.e. 对角态射 $\Delta_{X/S}$ 是正则浸入) 也意味着，作为增广到 $\mathcal{O}_X$ 的增广代数层，$\mathcal{P}_{X/S}^{\infty}$ 正则，i.e. $\varOmega_{X/S}^{1}$ 局部自由，且典范满同态
\[
S_{\mathcal{O}_X}\bigl(\varOmega_{X/S}^{1}\bigr) \longrightarrow \mathrm{gr}_{*}\bigl(\mathcal{P}_{X/S}^{\infty}\bigr)
\]
是同构，或最后，$X$ 的每个点都有一个开邻域，在其上增广代数层 $\mathcal{P}_{X/S}^{\infty}$ 同构于层 $\mathcal{O}_X[\![t_1, \dots, t_n]\!]$.

设 $s$ 是 $X$ 在 $S$ 上的截面，$\mathcal{J}$ 是它所定义的 $X$ 上的理想层 (为简化假设 $s(S)$ 闭)，则有 $\mathcal{O}_X$-增广代数层的典范同构：
\begin{equation}
s^{*}\bigl(\mathcal{P}_{X/S}^{n}\bigr) = \mathcal{O}_X/\mathcal{J}^{n+1}, \qquad s^{*}\bigl(\mathcal{P}_{X/S}^{\infty}\bigr) = \varprojlim_n \mathcal{O}_X/\mathcal{J}^{n+1}. \tag{4.4}
\label{eq:II.4.4}
\end{equation}
这些同构对基变换在显然的意义下是函子性的，并且 (考虑到这一点) 重新给出 $S$ 上代数层 $\mathcal{P}_{X/S}^{n}$ 的一个刻画. 例如若 $S = \mathrm{Spec}(k)$，$k$ 是域，则 $X$ 在 $S$ 上的截面 $s$ 的给出等价于 $X$ 在 $k$ 上有理点 $x$ 的给出，而前面的公式意味着有 $k$-代数同构
\begin{equation}
\mathcal{P}_{X/S}^{n}(x) = \mathcal{O}_x/\mathfrak{m}_x^{n+1}, \tag{4.5}
\label{eq:II.4.5}
\end{equation}
这证明``$X$ 相对 $S$ 的 $n$ 阶\emph{主部层} (faisceau des parties principales d'ordre n)\index[term]{zhubuceng@主部层}''这一名称之合理，该名称是给 $\mathcal{P}_{X/S}^{n}$ 的. 还由 \eqref{eq:II.4.4} 看到，若 $X$ 在 $s(S)$ 的所有点处微分光滑，则 $X$ 在 $s(S)$ 的所有点处光滑 (推论 \ref{cor:II.4.17})，逆命题也成立 (推论 \ref{cor:II.4.18}). 结合 \ref{cor:II.4.13}，容易由此推出：若 $X$ 是局部有限型 $S$-预概形，则 $X$ 在 $S$ 上光滑当且仅当它在 $S$ 上平坦且微分光滑. (N.B. 平坦性假设是本质的，取 $X$ 为 $S$ 的闭子预概形即可看出.)

再作为回忆指出，借助投影 $\mathrm{pr}_2 \colon X \times_S X \to X$ 在 $\mathcal{P}_{X/S}^{n}$ 上得到第二种代数结构，它由第一种结构经 $X \times_S X$ 的对称自同构诱导的环层 $\mathcal{P}_{X/S}^{n}$ 的典范对合 (involution) 得出. 用 $d_{X/S}^{n}$ 或简称 $d^n$ 记环层同态
\begin{equation}
d_{X/S}^{n} \colon \mathcal{O}_X \longrightarrow \mathcal{P}_{X/S}^{n}, \tag{4.6}
\label{eq:II.4.6}
\end{equation}
它对应于这第二种代数结构. 考虑到同构 \eqref{eq:II.4.4}，该同态把 $\mathcal{O}_X$ 的截面 $f$ 变成 $\mathcal{P}_{X/S}^{n}$ 的截面 $d^n(f)$，后者经 $X$ 在 $S$ 上的截面 $s$ 取原像后等同于 $f$ 在 $\Gamma(X, \mathcal{O}_X/\mathcal{J}^{n+1})$ 中的典范像. 这证明给 $d^n f$ 的``$f$ 的 $n$ 阶主部组 (système des parties principales d'ordre n)''名称之合理，特别是在 $S = \mathrm{Spec}(k)$ 的情形，如公式 \eqref{eq:II.4.5} 所考虑的.

最后注意到，同态 \eqref{eq:II.4.6} 可以看作 (相对常数预概形 $S$ 的) $\mathcal{O}_X$ 上的万有 $\leqslant n$ 阶\emph{微分算子} (opérateur différentiel d'ordre $\leqslant n$)\index[term]{weifensuanzi@微分算子}\footnote{关于本段的全部内容，可参阅 EGA IV 16.8 至 16.12.}，这里约定把 $\mathcal{O}_X$ 到模层 $F$ 的 $\leqslant n$ 阶微分算子称为可分解为
\[
D \colon \mathcal{O}_X \xrightarrow{d^n} \mathcal{P}_{X/S}^{n} \xrightarrow{u} F
\]
的层同态 $D$，其中 $u$ 是 $\mathcal{O}_X$-模层同态，且由 $D$ 唯一确定. 这个定义与直观的递归定义一致 ($D$ 是 $\leqslant n$ 阶微分算子，若对 $X$ 的开集 $U$ 上 $\mathcal{O}_X$ 的每个截面 $g$，$f \mapsto D(fg) - D(f)$ 是 $U$ 上的 $\leqslant n-1$ 阶微分算子). 由此推出，若 $X$ 在 $S$ 上微分光滑，则所有阶微分算子构成的环层具有微分流形 (differentiable manifold) 上微分演算中众所周知的简单结构，特别在局部有 $\mathcal{O}_X$-模层基，由可交换算子 $\delta/\delta x_i$ ($1 \leqslant i \leqslant n$) 的分次幂 (divided powers) 构成.

若 $S$ 是 $\mathbf{Q}$-代数层 ($\mathbf{Q}$ 为有理数域)，只需取诸 $\delta/\delta x_i$ 的通常多项式. 此时且仅在此时，极为特殊地，为使 $X$ 在 $S$ 上微分光滑，只需 $\varOmega_{X/S}^{1}$ 局部自由.
\end{remark}
\endgroup

\begin{remark}\label{rem:II.4.19}
本节引入的``正则浸入''、``正则理想''等术语遇到了相当激烈而普遍的反对 (Chevalley，Serre). 有人提议``Cohen-Macaulay 理想''、``Macaulay 理想''或``macaulay 理想''，这从道义上也将迫使我们采用``Cohen-Macaulay 浸入''或``Macaulay 浸入''. 然而这一术语与大全未来稿中已经使用的另一术语冲突，那里一个 (有限型) 态射称为在一点``Cohen-Macaulay''，若它在这一点平坦，且经过该点的纤维在该点的局部环是 Cohen-Macaulay 环. 在找到满意的方案之前，我们暂且保留本节引入的术语\footnote{这就是 EGA 0$_{\mathrm{IV}}$ 15.1.7 所采纳的.}.
\end{remark}

\section{基域的情形}

\begin{proposition}\label{prop:II.5.1}
设 $k$ 是域，$X$ 是 $k$ 上的有限型预概形，$x$ 是 $X$ 的点，$n$ 是 $X$ 在 $x$ 处的维数，$f \colon X \to \mathrm{Spec}\, k[t_1, \dots, t_n] = Y$ 是由 $\Gamma(X, \mathcal{O}_X)$ 的元素 $f_i$ 定义的态射. 下列条件等价 (且蕴含 $X$ 在 $k$ 上于 $x$ 处光滑，从而由 \ref{prop:II.3.1} 在 $x$ 处正则)：
\begin{enumerate}[(i)]
\item $f$ 在 $x$ 处平展.
\item $df_i$ 构成 $\varOmega_{X/k}^{1}$ 在 $x$ 处的基.
\item $df_i$ 在 $x$ 处生成 $\varOmega_{X/k}^{1}$.
\end{enumerate}
\end{proposition}

\begin{proof}
由于 (i) 蕴含 $X$ 在 $k$ 上于 $x$ 处光滑，(i)$\Rightarrow$(ii) 是 \ref{cor:II.4.8} 的特殊情形，(ii)$\Rightarrow$(iii) 平凡，剩下证明 (iii)$\Rightarrow$(i). 若 (iii) 成立，由引理 \ref{lem:II.4.1}，$f$ 在 $x$ 处非分歧，故 (把 $x$ 换成开邻域) 拟有限，从而由维数理由支配. 由于 $Y$ 正则，由 (\ref{thm:I.9.5} (ii)) 或 (\ref{cor:I.9.11}) 得出 $f$ 平展.
\end{proof}

\begin{corollary}\label{cor:II.5.2}
在 \ref{prop:II.5.1} 的预备条件下，假设 $k(x)$ 是 $k$ 的有限可分扩张，且 $f_i$ ($1 \leqslant i \leqslant n$) 定义 $\mathfrak{m}_x$ 的元素. 则前面的条件等价于
\begin{enumerate}[(i)]
\item[(iv)] $f_i$ 构成 $\mathfrak{m}_x$ 的生成元组 (或：$f_i$ 模 $\mathfrak{m}_x^2$ 构成 $k(x)$ 上的基 $\mathfrak{m}_x/\mathfrak{m}_x^2$).
\end{enumerate}
\end{corollary}

\begin{proof}
事实上，(iv)$\Rightarrow$(iii) 由正合序列
\begin{equation}
\mathfrak{m}_x/\mathfrak{m}_x^2 \longrightarrow \varOmega_{\mathcal{O}_x/k}^{1} \longrightarrow \varOmega_{k(x)/k}^{1} \longrightarrow 0 \tag{5.1}
\label{eq:II.5.1}
\end{equation}
得出，并考虑到 $k(x)$ 在 $k$ 上平展故 $\varOmega_{k(x)/k}^{1} = 0$. 另一方面 (ii) 蕴含 (iv)，因为 $X$ 和 $\mathrm{Spec}\bigl(k(x)\bigr)$ 都在 $k$ 上于 $x$ 处光滑，由 \ref{thm:II.4.10} (iv) 可以在上述正合序列左边加一个 0.
\end{proof}

\begin{corollary}\label{cor:II.5.3}
设 $x$ 是 (在 $k$ 上有限型的) $X$ 的点. 若 $X$ 在 $k$ 上于 $x$ 处光滑，则 $\mathcal{O}_x$ 正则，逆命题在 $k(x)$ 是 $k$ 的有限可分扩张时成立.
\end{corollary}

\begin{proof}
事实上，逆命题由 \ref{cor:II.5.2} 得出，取 $\mathfrak{m}_x$ 的正则生成元组 ($f_i$) (N.B. 除 \ref{cor:II.5.2} 外，也可以援引定理 \ref{thm:II.4.15}). 于是：
\end{proof}

\begin{proposition}\label{prop:II.5.4}
设 $X$ 是 $k$ 上的有限型预概形. 若 $X$ 在 $k$ 上光滑，则它正则，逆命题在 $k$ 完美 (perfect) 时成立.
\end{proposition}

\begin{proof}
对逆命题，注意由 \ref{cor:II.5.3}，$X$ 在 $k$ 上于每个闭点光滑，从而处处光滑，因为光滑的点集是开的.
\end{proof}

\begin{theorem}\label{thm:II.5.5}
设 $X$ 是 $k$ 上的有限型预概形，$x$ 是 $X$ 的点，$n$ 是 $X$ 在 $x$ 处的维数，$k'$ 是 $k$ 的完美扩张. 下列条件等价：
\begin{enumerate}[(i)]
\item $X$ 在 $k$ 上于 $x$ 处光滑.
\item $\varOmega_{X/k}^{1}$ 在 $x$ 处是秩为 $n$ 的自由模.
\item[(ii bis)] $\varOmega_{X/k}^{1}$ 在 $x$ 处由 $n$ 个元素生成.
\item $X$ 在 $k$ 上于 $x$ 处微分光滑.
\item 存在 $x$ 的开邻域 $U$，使得 $U \otimes_k k'$ 正则 (i.e. 其各点的局部环正则).
\end{enumerate}
\end{theorem}

\begin{proof}
由 \ref{thm:II.4.3} (ii) 有 (i)$\Rightarrow$(ii)，(ii)$\Rightarrow$(ii bis) 平凡，(ii bis)$\Rightarrow$(i) 由 \ref{prop:II.5.1} 得出. 由于 $X$ 在 $k$ 上平坦，由 \ref{cor:II.4.18} 有 (i)$\Leftrightarrow$(iii). 有 (i)$\Rightarrow$(iv)，因为光滑对基扩张不变且蕴含正则，而 (iv)$\Rightarrow$(i) 因为由命题 \ref{prop:II.5.4}，$U \otimes_k k'$ 在 $k'$ 上简单，故由 \ref{cor:II.4.13}，$U$ 在 $k$ 上简单.
\end{proof}

对 $x$ 取 $X$ 的生成点 (假设 $X$ 不可约)，得到：

\begin{corollary}\label{cor:II.5.6}
设 $K$ 是域 $k$ 上有限型代数的局部化所成的 Artin 局部环 (例如 $K$ 可以是 $k$ 的有限型扩张)，设 $n$ 是其剩余域在 $K$ 上的超越次数. 下列条件等价：
\begin{enumerate}[(i)]
\item $K$ 是 $k$ 的纯超越扩张 (purely transcendental extension) $k(t_1, \dots, t_n)$ 的有限可分扩张.
\item $\varOmega_{K/k}^{1}$ 是秩为 $n$ 的自由 $K$-模.
\item[(ii bis)] $\varOmega_{K/k}^{1}$ 是可由 $n$ 个元素生成的 $K$-模.
\item $K \otimes_k K$ 按增广理想 $K \otimes_k K \to K$ 的幂所定义的拓扑的完备化 $\mathcal{O}'$ 是``正则''增广 $K$-代数，i.e. 同构于 $K$ 上的形式幂级数代数 (N.B. 若 $K$ 是域，这等价于说 $\mathcal{O}'$ 是正则局部环).
\item $K$ 是 $k$ 的可分扩张.
\end{enumerate}
\end{corollary}

\begin{proof}
事实上，总可以把 $K$ 看作 $k$ 上不可约有限型概形 $X$ 的生成点的局部环，而所考虑的条件正是 \ref{thm:II.5.5} 中同名条件，在 (iv) 中对 $k'$ 取 $k$ 的代数闭包. 只有``$K$ 在 $k$ 上可分 $\Rightarrow$ $X$ 在 $k$ 上于 $x$ 处光滑''这一蕴含需要证明. 借助 \ref{cor:II.4.13} 立即归结到基域为 $k'$ 从而代数闭的情形，此时存在 $X$ 在 $k$ 上有理点 $a$. 但此时由 \ref{prop:II.5.4}，$X$ 在 $k$ 上于 $a$ 处光滑，更不用说它在 $k$ 上于生成点 $x$ 处光滑，证毕\footnote{注意对 $k$ 所作的假设是不可缺少的，如下例所示：$k = \mathbf{R}$，$X = \mathrm{Spec}(k)$，$Y = \mathrm{Spec}\bigl(k[t]/(t^2+1)\bigr)$. 则 $Y$ 是 $X$ 的平展覆盖但不平凡，尽管 $k$ 在 $k(t)$ 中代数闭.}.
\end{proof}

注意，当 $K$ 是 $k$ 的有限型扩张时，(i)，(ii)，(ii bis)，(iv) 的等价性是众所周知的，但我们没有使用这些已知等价性中的任何一个. 当然，命题 \ref{prop:II.5.1} 包含作为特殊情形：元素组 $x_i$ ($1 \leqslant i \leqslant n$) 是 $K$ 在 $k$ 上的``可分超越基'' (separating transcendence base)，当且仅当 $dx_i$ 构成 $K$-模 $\varOmega_{K/k}^{1}$ 的基，这事实别处也都知道.

\begin{corollary}\label{cor:II.5.7}
设 $X$ 是域 $k$ 上的有限型预概形. 为使 $X$ 在 $k$ 上光滑，必要且充分条件是 $\varOmega_{X/k}^{1}$ 局部自由，且 $X$ 的不可约分支的生成点的局部环是 $k$ 的可分扩张 (最后一个条件在 $k$ 完美且 $X$ 约化时自动满足).
\end{corollary}

\begin{proof}
可以假设 $X$ 连通，设 $n$ 是假设局部自由的 $\varOmega_{X/k}^{1}$ 的秩. 由假设和 \ref{cor:II.5.6}，它也是 $X$ 的生成点的局部环所定义的 $k$ 的扩张的超越次数，故 $X$ 的所有不可约分支维数为 $n$. 于是由 \ref{thm:II.5.5} 得出结论.
\end{proof}

注意，若 $K$ 是 $k$ 的有限 (不必可分的) 扩张，则 $\varOmega_{K/k}^{1}$ 是自由 $k$-模，故令 $X = \mathrm{Spec}(K)$，$\varOmega_{X/k}^{1}$ 是局部自由层，且 $X$ 约化，而 $X$ 未必在 $k$ 上光滑. 再把标量扩张到 $k$ 的代数闭包，得到类似的例子，其中 $k$ 代数闭，但 $X$ 却不再约化.

\begin{corollary}\label{cor:II.5.8}
设 $X$ 是域 $k$ 上的有限型预概形，$x$ 是 $X$ 的点，$n$ 是 $X$ 在 $x$ 处的维数，$p$ 是 $\mathcal{O}_x$ 的维数，i.e. $x$ 在 $X$ 中的闭包 $Y$ 的余维数；故 $n - p$ 是 $k(x)$ 在 $k$ 上的超越次数. 设 $f_i$ ($1 \leqslant i \leqslant n$) 是 $\mathcal{O}_x$ 的元素，对 $1 \leqslant i \leqslant p$ 有 $f_i \in \mathfrak{m}_x$. 下列条件等价：
\begin{enumerate}[(i)]
\item 由诸 $f_i$ 定义的 $x$ 处的态射芽
\[
X \longrightarrow \mathrm{Spec}\bigl(k[t_1, \dots, t_n]\bigr)
\]
在 $x$ 处平展.
\item $f_i$ ($1 \leqslant i \leqslant p$) 生成 $\mathfrak{m}_x$，i.e. 构成 $\mathcal{O}_x$ 的正则参数系 (regular system of parameters)，且 $f_j$ ($p+1 \leqslant j \leqslant n$) 在 $k(x)$ 中的类构成可分超越基 (i.e. $d\overline{f}_j$ ($p+1 \leqslant j \leqslant n$) 构成 $\varOmega_{k(x)/k}^{1}$ 的基，或生成 $\varOmega_{k(x)/k}^{1}$).
\end{enumerate}
\end{corollary}

\begin{proof}
假设 (i) 成立. 由此得出 $df_i(x)$ 构成 $\varOmega_{X/k}^{1}(x)$ 的基 (\cref{cor:II.4.8})，故它们在 $\varOmega_{k(x)/k}^{1}$ 中的像 $d\overline{f}_i(x)$ 生成该 $k$-向量空间. 由于对 $1 \leqslant i \leqslant p$，$\overline{f}_i$ 为零，只需取 $p+1 \leqslant i \leqslant n$ 的 $d\overline{f}_i(x)$. 由于 $k(x)$ 在 $k$ 上的超越次数为 $n - p$，由推论 \ref{cor:II.5.8} 判则 (iii) (应用于 $K = k(x)$) 得出 $Y$ 在 $k$ 上于其生成点 $x$ 处光滑，且 $d\overline{f}_i(x)$ ($p+1 \leqslant i \leqslant n$) 构成 $\varOmega_{k(x)/k}^{1}$ 在 $k(x)$ 上的基. 于是 \ref{prop:II.4.9} 的条件 (ii) 成立，从而条件 (iii) 也成立，特别地 $f_i$ ($1 \leqslant i \leqslant p$) 构成 $\mathfrak{m}_x$ 的生成元组. 由于 $\mathcal{O}_x$ 的维数为 $p$，它们构成 $x$ 处的正则参数系. 这证明 (ii).
\end{proof}

假设 (ii) 成立. 由正合序列 \eqref{eq:II.5.1}，由此得出 $df_i(x)$ 生成 $\varOmega_{X/k}^{1}$，从而由命题 \ref{prop:II.5.1} 得 (i).

\begin{corollary}\label{cor:II.5.9}
设 $X$ 是域 $k$ 上的有限型预概形，$x$ 是 $X$ 的点，$n$ 是 $X$ 在 $x$ 处的维数，$p$ 是 $\mathcal{O}_x$ 的维数，i.e. $x$ 在 $X$ 中的闭包 $Y$ 的余维数，故 $n - p$ 是 $k(x)$ 在 $k$ 上的超越次数. 下列条件等价：
\begin{enumerate}[(i)]
\item $\mathcal{O}_x$ 正则且 $k(x)$ 是 $k$ 的可分扩张.
\item $X$ 在 $k$ 上于 $x$ 处光滑，且典范同态
\[
\mathfrak{m}_x/\mathfrak{m}_x^2 \longrightarrow \varOmega_{\mathcal{O}_x/k}^{1} \otimes_{\mathcal{O}_x} k(x) = \varOmega_{X/k}^{1}(x)
\]
单射.
\item 存在 $f_i \in \mathcal{O}_x$ ($1 \leqslant i \leqslant n$)，对 $1 \leqslant i \leqslant p$ 有 $f_i \in \mathfrak{m}_x$，使得由诸 $f_i$ 定义的 $X$ 到 $\mathrm{Spec}\bigl(k[t_1, \dots, t_n]\bigr)$ 的 $x$ 处态射芽在 $x$ 处平展 (i.e. 由 \ref{prop:II.5.1}，使得 $df_i(x)$ 生成 $\varOmega_{X/k}^{1}(x)$).
\item 存在 $f_i \in \mathcal{O}_x$ ($1 \leqslant i \leqslant n$)，使得 $f_i$ ($1 \leqslant i \leqslant p$) 生成 $\mathfrak{m}_x$，且 $df_j(x)$ ($p+1 \leqslant j \leqslant n$) 在 $k(x)$ 上生成 $\varOmega_{k(x)/k}^{1}$.
\end{enumerate}
\end{corollary}

\begin{proof}
(iii) 与 (iv) 的等价性由推论 \ref{cor:II.5.8} 得出，且由 \ref{prop:II.4.9}，这些条件还等价于``$X$ 在 $k$ 上于 $x$ 处光滑且 \ref{thm:II.4.10} 的条件 (ii) 成立''. 故它们等价于``$X$ 在 $k$ 上于 $x$ 处光滑且 \ref{thm:II.4.10} 的条件 (iv) 成立''，从而等价于 \ref{cor:II.4.8} (ii). 或等价于``$X$ 在 $k$ 上于 $x$ 处光滑且 \ref{thm:II.4.10} 的条件 (i) 成立''，在这里它简单意味着 $k(x)$ 在 $k$ 上可分. 这蕴含 \ref{cor:II.5.9} (i)，剩下证明 \ref{cor:II.5.9} (i) 蕴含它，i.e. 证明：
\end{proof}

\begin{corollary}\label{cor:II.5.10}
设 $x$ 是域 $k$ 上有限型预概形的点，使得 $k(x)$ 在 $k$ 上可分. 为使 $X$ 在 $k$ 上于 $x$ 处光滑，必要且充分条件是它在 $x$ 处正则 (i.e. 局部环 $\mathcal{O}_x$ 正则).
\end{corollary}

\begin{proof}
事实上，若如此，显然可以找到 $f_i \in \mathcal{O}_x$ ($1 \leqslant i \leqslant n$) 满足条件 \ref{cor:II.5.9} (iv).
\end{proof}

\begin{remarkstar}
\textbf{勘误 (Errata).} 在本节 \ref{cor:II.5.6} 的证明中，我们使用了``代数闭域上非空的有限型约化概形至少有一个正则 (从而光滑) 点''这一事实，该事实通常用微分方法证明 (借助 Zariski 的``$X$ 的正则点集是开的''这一定理). 若要避免循环论证，必须证明：若 $K/k$ 是有限型可分扩张，且 $f_i \in K$ 使得 $d_{K/k} f_i$ ($1 \leqslant i \leqslant n$) 构成 $\varOmega_{K/k}^{1}$ 的基，则 $n$ 是 $K$ 在 $k$ 上的超越次数，i.e. 诸 $f_i$ 代数无关 (algebraically independent). 用 Mac Lane 判则证明这一事实是众所周知的，参见 Bourbaki，Algèbre，Chap. V，§9，th. 2：取次数最小的多项式 $g \in k[t_1, \dots, t_n]$ 使得 $g(f_1, \dots, f_n) = 0$，于是有
\[
\sum_i \frac{dg}{dt_i}(f_1, \dots, f_n)\, df_i = 0,
\]
从而 (因为 $df_i$ 构成 $\varOmega_{K/k}^{1}$ 的基) $dg/dt_i$ 在 $(f_1, \dots, f_n)$ 处取值为零，故由 $g$ 的极小性它们为零，于是若 $k$ 特征为 0 则 $g = 0$，若 $k$ 特征 $p \neq 0$ 则 $g = h(t_1^p, \dots, t_n^p)$. 用 Mac Lane 判则，多项式 $h \in k[t_1, \dots, t_n]$ 也零化 $(f_1, \dots, f_n)$，从而仍由 $g$ 的极小性得 $g = 0$.
\end{remarkstar}

